\scanpage{226}
\mspage{176}
% L1–3
à ces indications, quitte à faire des vérifications circonstanciées par la suite. Au besoin, j'utiliserai dès maintenant que pour $i\in\mathbb N$, $i\geq1$,
% L4
\[
\pi_{\unsure{i|n}}(\deleted{\unsure{X}}f,x):\pi_i(X,x)\unsure{-|\to}\pi_i(Y,f(x))
\]
% L5–6
\unsure{peuvent|peut} \unsure{s'interpréter|s'interprète} comme \inserted{\unsure{proviennent de|provient de}} un hom. de systèmes locaux
% L7
\[
\underset{\sim}{\pi}_i(X,-)\longrightarrow f^*(\underset{\sim}{\pi}_i(Y,-))
\]
% L8–13
sur $X$. La condition a) \inserted{du th. 1} s'interprète en disant que cet hom est un \uline{iso} pour $i\in[1,n]$, et la condition a$'$) se borne alors à rappeler qu'il suffit de tester qu'il en est ainsi pour \uline{un} point de chaque composante connexe de $X$.

% L14–16
L'énoncé du th. 1 garde un sens pour $n=1$, et dans ce cas a été ``rappelé'' précédemment (et est $\pm$ tautologique).
% L17–24
Cela\unsure{$=$} \unsure{va} nous permettre de procéder par récurrence sur $n$, en supposant que le th. est prouvé pour $n-1$. Cela signifie que l'on peut supposer que les conditions a) \uline{et} b) sont satisfaites \inserted{par $f$} pour $i\in[1,n-1]$, et se borner à prouver que \inserted{sous ces conditions} \deleted{\unsure{a}} la \unsure{c.} a) pour $i=n$, équivaut à la condition b) pour $\unsure{j|i}=n$.
% archive sheet number: 226
